\documentclass[10pt,a4paper]{amsart}
\usepackage[margin=19mm]{geometry}
\usepackage{amsmath,amssymb,mathtools}
\usepackage[scr=rsfs,cal=euler]{mathalfa}
\usepackage{xfrac}
\usepackage[unicode,hidelinks,bookmarks=false]{hyperref}
\newcommand{\ie}{i.e.\ }
\newcommand{\cf}{cf.\ }
\newcommand{\resp}{respectively\ }
\newcommand{\Subsection}{\subsection}
\providecommand{\nobreakdash}{\nobreak\hbox{-}}
\setlength{\emergencystretch}{2em}
\multlinegap0pt
\usepackage{bm}
\usepackage{bbm}
\usepackage{dsfont}
\usepackage{xspace}
\usepackage{cancel}
\usepackage{mathtools}
\usepackage{amsthm}
\makeatletter
\@ifundefined{thm}{\newtheorem{thm}{Theorem}}{}
\@ifundefined{lem}{\newtheorem{lem}[thm]{Lemma}}{}
\@ifundefined{prop}{\newtheorem{prop}[thm]{Proposition}}{}
\makeatother
\title[Lower bounds for fractional Orlicz-type eigenvalues]{Lower bounds for fractional Orlicz-type eigenvalues}
\author{Ariel Salort}
\date{Source: arXiv:2502.15423v2 (CC BY 4.0)}
\begin{document}
\maketitle

\noindent\emph{Source and licence.} Lower bounds for fractional Orlicz-type eigenvalues. Version arXiv:2502.15423v2 (CC BY 4.0), distributed under the Creative Commons Attribution 4.0 International licence, as recorded in the arXiv licence metadata for this exact version and shown on the abstract page https://arxiv.org/abs/2502.15423v2. Full rights notice: licenses/arxiv-2502.15423-CC-BY-4.0.txt.\par\smallskip

\noindent\textit{Editorial scope.} Counted unit: the Thm whose source label is \texttt{teo2} in the article (source0.tex lines 963--982, source label \texttt{teo2}), delivered together with the article's own complete original proof of it (source0.tex lines 984--1023, one \texttt{proof} environment of 40 lines). No mathematics is written, repaired or re-derived: statement and proof are the article's own text line for line. Required context retained uncounted: minimi.intro (lines 159--161); teo1.intro (lines 194--212); cond1 (lines 579--581); func.e (lines 583--585); lema1 (lines 593--616); morrey (lines 619--625); minimi (lines 740--742); eq (lines 746--748); desig0 (lines 890--899). Every definition, display and label of those lines is retained unchanged. Omitted from this package and not redistributed: 1--158, 162--193, 213--578, 582--582, 586--592, 617--618, 626--739, 743--745, 749--889, 900--962, 983--983, 1024--1243. Nothing inside a retained excerpt is omitted or altered: every retained body is delivered exactly as the article's lines are written, so no omission receipt of any kind accompanies this package. Citations: the retained excerpts carry no citation command at all, so no bibliography entry of the article is transcribed and no reference is redistributed. Reference adaptation: none was needed and none was made; every reference inside the delivered unit resolves inside it, so no unresolved reference or citation is produced. Printed slips kept literal: no printed word, symbol, hypothesis or number of the counted statement or of its proof was altered. Layout and class: the article's own class is \texttt{amsart} with packages including lineno, inputenc, graphicx, hyperref, geometry; the wrapper is the lane's pinned \texttt{amsart} editorial wrapper at 10pt on A4, which restates the theorem-like environments the retained text needs and adds the article's own macro definitions that the retained text needs (none), with extra wrapper packages (bm, bbm, dsfont, xspace, cancel, mathtools). The delivered numbering is the unit's own. Assets: no retained span contains a figure environment, a graphics-inclusion command, a PDF-page inclusion, a table or a TikZ picture; no image file is copied into this package, the package declares no asset dependency and no image file is redistributed with it. Licence: version arXiv:2502.15423v2 of this article is distributed under the Creative Commons Attribution 4.0 International licence, as recorded in the arXiv licence metadata for this exact version and shown on the abstract page https://arxiv.org/abs/2502.15423v2. A full rights notice is delivered at licenses/arxiv-2502.15423-CC-BY-4.0.txt. Characters: every retained excerpt is pure ASCII, so the delivered engine, XeLaTeX with its default Latin Modern text font, reports no missing character.
\begin{equation} \label{minimi.intro}
\lambda_{1,\alpha}^s = \inf\left\{\frac{1}{\alpha}\iint_{\mathbb{R}^{2n}} A(|D^s  u|)\, d\nu_n \colon u\in C_c^\infty(\Omega), \int_\Omega  \omega A(|u|)\, dx = \alpha\right\},
\end{equation}

\begin{thm}  \label{teo1.intro}
Let $s\in (0,1)$ and let $A$ be a Young function satisfying \eqref{cond1}. Let $\Omega\subset \mathbb{R}^n$ be an open bounded domain with inner radius $r_\Omega$. Given  $\omega\in L^1(\mathbb{R}^n)$ and $\alpha>0$, consider the critical value $\lambda_{1,\alpha}^s$ defined in \eqref{minimi.intro}. The following holds:
\begin{itemize}
\item[i)] There exists a unique $\alpha_0>0$ satisfying the equation
$$
\alpha_0 \lambda_{1,\alpha_0}^s = r_\Omega^n.
$$

\item[ii)] 
Assume that $i_0(A)>\frac{n}{s}$ when $\alpha\leq \alpha_0$ or $i_\infty(A) >\frac{n}{s}$ when $\alpha> \alpha_0$. Then, there exists a positive, computable constant $C=C(n,s,A)$ such that
\begin{equation*} %\label{ineq.intro}
\frac{C}{\|\omega\|_{L^1(\Omega)}} \frac{r_\Omega^n}{M_A(r_\Omega^s)} \le \lambda_{1,\alpha}^s.
\end{equation*}


\noindent In particular, this holds when $i_0(A)>\frac{n}{s}$ if $\alpha\ll1$ or when $i_\infty(A) >\frac{n}{s}$ if $\alpha\gg1$.
\end{itemize}

\end{thm}

\begin{equation} \label{cond1}
\int^\infty \left(\frac{t}{A(t)}\right)^\frac{s}{n-s}\,dt<\infty
\end{equation}

\begin{equation} \label{func.e}
E(t)=t^\frac{n}{n-s}\int_t^\infty \frac{\tilde A(\tau)}{\tau^{1+\frac{n}{n-s}}}\,d\tau\quad \text{ for } t\geq 0.
\end{equation}

\begin{lem} \label{lema1}
Let $s\in (0,1)$ and let $A$ be a Young function. Assume \eqref{cond1}. Then

\begin{itemize}
  \setlength\itemsep{0.5em}
\item[(i)] The function $\Psi_s$ is non-decreasing.

\item[(ii)] Define the Young function $B=\tilde E$. Then
$
\Psi_s(r)\approx r^s B^{-1}(r^{-n}) \text{ for } r>0.
$

\item[(iii)] If $i_\infty(A)>\frac{n}{s}$ then $B\simeq A$ near infinity, and
$
\Psi_s(r) \approx  r^s A^{-1}(r^{-n})  \text{ for } 0<r\leq 1.
$

\item[(iv)] If $i_0(A)>\frac{n}{s}$ then $B\simeq A$ near $0$, and
$
\Psi_s(r) \approx  r^s A^{-1}(r^{-n}) \text{ for } r\geq 1.
$ 
\end{itemize}
 
\end{lem}

\begin{prop} \label{morrey}
Let $s\in (0,1)$ and let $A$ be a Young function satisfying  \eqref{cond1}. Then, $W^s L^A(\mathbb{R}^n)\subset C^{\Psi_s(\cdot)}(\mathbb{R}^n)$. Moreover,  for any $u\in W^s L^A(\mathbb{R}^n)$ and $x,y\in \mathbb{R}^n$ it holds that
$$
|u(x)-u(y)|\leq C_M |x-y|^s B^{-1}\left( \frac{1}{|x-y|^n} \iint_{\mathbb{R}^{2n}} A(D^s u(z,w))\,d\nu_n(z,w) \right)
$$
for some constant $C_M$ depending on $n$ and $s$, where the Young function $B$ is given by $B(t)=\tilde E(t)$, being $E$ the Young function defined in \eqref{func.e}.
\end{prop}

\begin{equation} \label{minimi}
\lambda_{1,\alpha}^s = \inf\left\{\frac{1}{\alpha}\iint_{\mathbb{R}^{2n}} A(|D^s  u|)\, d\nu_n \colon u\in C_c^\infty(\Omega), \int_\Omega  \omega A(|u|)\, dx = \alpha\right\}.
\end{equation}

\begin{equation} \label{eq}
\iint_{\mathbb{R}^{2n}} A(|D^s u_\alpha^s|)\,d\nu_n = \lambda_{1,\alpha}^s \int_\Omega \omega A(|u_\alpha^s|)\,dx.
\end{equation}

\begin{align}\label{desig0}
\begin{split}
|u_\alpha(x_0)|&\leq C_M |x_0-y|^s B^{-1}\left(  \frac{\lambda_{1,\alpha} }{|x_0-y|^n}  \int_\Omega \omega A(|u_\alpha|)\,dx\right)\\
&= C_M |x_0-y|^s B^{-1}\left(   \frac{ \alpha \lambda_{1,\alpha}}{|x_0-y|^n}\right)\\
&=
C_M (\alpha\lambda_{1,\alpha})^\frac{s}{n} \left((\alpha \lambda_{1,\alpha})^{-\frac{1}{n}} |x_0-y| \right)^s B^{-1}\left(  \left( (\alpha \lambda_{1,\alpha}^{-\frac1n} |x_0-y| \right)^{-n} \right)\\
&\leq  
c_1  C_M \left(\alpha \lambda_{1,\alpha}  \right)^{\frac{s}{n}}\Psi_s \left( |x_0-y| \left( \alpha \lambda_{1,\alpha}^s   \right)^{-\frac{1}{n}}\right).
\end{split}
\end{align}

\begin{thm}  \label{teo2}
Let $s\in (0,1)$, $\alpha>0$ and let $A$ be a Young function satisfying \eqref{cond1}. Given $\omega \in L^1(\Omega)$ consider the critical value $\lambda_{1,\alpha}^s$  defined in \eqref{minimi}. It holds that
\begin{itemize}
\item[i)] There exists a unique $\alpha_0>0$ satisfying the equation
$$
\alpha_0 \lambda_{1,\alpha_0}^s = r_\Omega^n.
$$

\item[ii)] 
Assume that $i_0(A)>\frac{n}{s}$ when $\alpha\leq \alpha_0$, or $i_\infty(A) >\frac{n}{s}$ when $\alpha> \alpha_0$.  Then, there exists a  $C>0$ depending only on  $s$, $n$ and $A$ such that
\begin{equation} \label{ineq.2}
\frac{r_\Omega^n}{\alpha} A\left( \frac{1}{C r_\Omega ^s}A^{-1}\left( \frac{\alpha}{ \|\omega\|_{L^1(\Omega)}} \right) \right)
\leq 
    \lambda_{1,\alpha}^s.
\end{equation}


\noindent In particular, this holds when $i_0(A)>\frac{n}{s}$ if $\alpha\ll1$ and when $i_\infty(A) >\frac{n}{s}$ if $\alpha\gg1$.
\end{itemize}
\end{thm}

\begin{proof}
Fixed $\alpha>0$, let $u_\alpha^s \in W^s_0 L^A(\Omega)$ be a minimizer of \eqref{minimi} such that $\int_\Omega \omega A(|u_\alpha^s|)\,dx =\alpha$, i.e., the pair $(u_\alpha^s, \lambda_{1,\alpha}^s)$ satisfies equation \eqref{eq}, where  $\lambda_{1,\alpha}^s$ is defined in \eqref{minimi}. Since $s\in(0,1)$ is fixed, for simplicity we will drop the dependence in $s$.

\noindent In light of Proposition \ref{morrey}  $u_\alpha$ is  continuous and hence there exists $x_0\in \Omega$ such that
$$
|u_\alpha(x_0)|=\max\{|u_\alpha(x)|\colon x \in \mathbb{R}^n\}>0.
$$
Since
$$
\alpha = \int_\Omega \omega A(|u_\alpha|)\,dx \leq A(|u_\alpha(x_0)|) \|\omega\|_{L^1(\Omega)},
$$
using expression \eqref{desig0} and the fact the $A$ is no decreasing,  we get that
\begin{align*}
\alpha &\leq
A\left( c_1  C_M \left(\alpha \lambda_{1,\alpha}  \right)^{\frac{s}{n}}\Psi_s \left( |x_0-y| \left( \alpha \lambda_{1,\alpha}   \right)^{-\frac{1}{n}}\right) \right) \|\omega\|_{L^1(\Omega)}
\end{align*}
for any $y\in \partial\Omega$, and therefore,  denoting by $r_\Omega$ the inner radius of $\Omega$, we get
\begin{equation} \label{desig11}
\alpha \|\omega\|_{L^1(\Omega)}^{-1}  
\leq
A\left( c_1  C_M \left(\alpha \lambda_{1,\alpha} \right)^{\frac{s}{n}}\Psi_s \left( r_\Omega \left( \alpha \lambda_{1,\alpha}   \right)^{-\frac{1}{n}}\right) \right).
\end{equation}
As in the proof of Theorem \ref{teo1.intro},  there exists $\alpha_0>0$ such that  $r_\Omega^n  \leq \alpha \lambda_{1,\alpha}$ when $\alpha>\alpha_0$ and $r_\Omega^n \geq \alpha \lambda_{1,\alpha}$ when $\alpha<\alpha_0$.

\medskip

{\bf Case }$\alpha>\alpha_0$. Assuming  $i_\infty(A) >\frac{n}{s}$, by  Lemma \ref{lema1} (iii) we get that $\Psi_s(t) \approx t^s A^{-1} (t^{-n})$ for any $t>0$. Then, there exists $c_2=c_2(n,s)>0$ for which  \eqref{desig11} yields
\begin{align*}
\alpha \|\omega\|_{L^1(\Omega)}^{-1}\leq 
A\left( C r_\Omega^s A^{-1} \left(   r_\Omega^{-n}\alpha \lambda_{1,\alpha} \right) \right)
\end{align*}
where $C=c_1c_2C_M$. Since $A$ is non decreasing, the previous expression gives that 
\begin{align*}
\frac{r_\Omega^n}{\alpha} A\left( \frac{1}{C r_\Omega ^s}A^{-1}\left( \frac{\alpha}{ \|\omega\|_{L^1(\Omega)}} \right) \right)
\leq 
    \lambda_{1,\alpha}.
\end{align*} 

{\bf Case }$\alpha>\alpha_0$. Assuming that $i_0(A) >\frac{n}{s}$, the bound follows analogously.
\end{proof}

\end{document}
